\section{MIDI Domain Logic and Encoding}\label{sec:implementation-music-logic}

A primary objective of the implementation is to bridge the gap between high-level musical notation and the
low-level binary encoding required by the MIDI standard.
This section details the mathematical formulas and mapping strategies implemented in the \texttt{dsl\_music}
module to ensure that the compiler's output is musically accurate.

\subsection{Note to MIDI Pitch Mapping}

The DSL allows composers to specify pitches using traditional letter names (A--G), accidentals (\#, b), and
octave numbers (0--8).
The compiler converts these literals into a single numeric byte (0--127) using a deterministic formula.

The implementation defines the \texttt{music::Note} structure as a combination of three fields:
\begin{itemize}
  \item \textbf{\texttt{Pitch} enum:} Stores the semitone offset within an octave (\texttt{C=0},
    \texttt{D=2}, \texttt{E=4}, \texttt{F=5}, \texttt{G=7}, \texttt{A=9}, \texttt{B=11}).
  \item \textbf{\texttt{Accidental} enum:} Stores a numeric signed offset (\texttt{Flat=-1},
    \texttt{Natural=0}, \texttt{Sharp=1}).
  \item \textbf{\texttt{Octave} integer:} Stores the numeric octave.
\end{itemize}

The mapping formula used in \texttt{Note::midi\_number()} is:
\begin{equation}
  P = 12 \times (O + 2) + V_p + V_a
\end{equation}
Where $P$ is the final MIDI pitch, $O$ is the octave, $V_p$ is the pitch-class value, and $V_a$ is the accidental value.
The constant offset of \textbf{+2} is applied to the octave to align with the standard MIDI convention where
\texttt{C4} is assigned the value 72 (a convention popularized by software such as Ableton Live, meaning
"Middle C" or \texttt{C3} is 60). By centralizing this calculation at the library level, the compiler ensures
that all musical logic throughout the pipeline is absolutely consistent.

\subsection{Instrument and Drum Mapping}

The language's \texttt{using} keyword allows the composer to specify the instrumental identity of a track.
This intent is mapped directly to the General MIDI (GM) Level 1 specification.

The implementation utilizes two primary mapping strategies:
\begin{enumerate}
  \item \textbf{Melodic Instruments:} Standard instruments are mapped to their GM program numbers. For
    example, \texttt{Piano} maps to 1, \texttt{Guitar} to 24, \texttt{Bass} to 32, and \texttt{Violin} to 40.
    These mappings are stored in the \texttt{music::Instrument} enum. During backend serialization, they are
    emitted as \textit{Program Change} events on standard MIDI channels (0--8, 11--15).
  \item \textbf{Percussion Logic:} Unlike melodic tracks, the MIDI standard rigidly reserves \textbf{Channel
    10} specifically for unpitched percussion. If a track is defined as \texttt{using drums}, the backend
    automatically routes all its events to Channel 10 and skips emitting a Program Change event. Furthermore,
    individual drum literals in the DSL are mapped to specific GM "Note-On" pitches as defined by the standard drum map:
    \begin{itemize}
      \item \texttt{Kick} $\rightarrow$ 36
      \item \texttt{Snare} $\rightarrow$ 38
      \item \texttt{Hihat} $\rightarrow$ 42
      \item \texttt{Crash} $\rightarrow$ 49
      \item \texttt{Ride} $\rightarrow$ 51
    \end{itemize}
\end{enumerate}
This explicit mapping guarantees that the compiled MIDI artifact will sound correct regardless of the Digital
Audio Workstation (DAW) or synthesizer used to play it back.
